\section{A differential inequality for the quadratic model}
\label{sec:model}

The transport law of Lemma~\ref{lem:fit-transport} contains the
coefficient
\[
 C(\lambda,\kappa,r)=\frac{R}{G},\qquad
 D\kappa=-C(\lambda,\kappa,r)\bigl(\lambda-\phi'(h)\bigr).
\]
When this identity is differentiated with respect to width, the
implicit changes of the fitted parameters contribute terms proportional
to $\kappa_r$. The remaining source contains the explicit model
derivative $C_r$. The shape comparison in Section~\ref{sec:transport}
will use its sign, which is the purpose of this section.

From this point onward $\lambda$ and $\kappa$ are independent model
parameters. The model quantities are
\[
 H=H(\lambda,\kappa,r),\qquad
 P=H_\lambda,\quad Q=H_\kappa,\quad
 R=P_r,\quad S=Q_r,\quad G=PS-QR,
\]
and $P,Q,R,S,G>0$ by Theorem~\ref{thm:quadratic-fit}.
Throughout this section, an $r$-subscript on a model function denotes a partial derivative
with $\lambda,\kappa$ fixed. Derivatives along a characteristic
will be identified explicitly.

\begin{theorem}
\label{thm:model-inequality}
For every \(\lambda,\kappa\in\mathbb R\) and \(r>0\),
\begin{equation}
 \left(\frac{R}{G}\right)_r\leq 0.
 \label{model:inequality}
\end{equation}
\end{theorem}

We prove the inequality by transporting the sign of
\(\mathcal W=RS_r-SR_r\) from small width. Endpoint identities and
a Riccati linearization determine the sign at any possible zero;
a characteristic argument then excludes such zeros when
\(\kappa\ne0\), and continuity treats \(\kappa=0\).

\subsection{Transport and the quantity whose sign is required}

Set
\begin{equation}
 \mathcal W=RS_r-SR_r,\qquad
 \alpha=\frac{r}{r^2-H^2},\qquad
 \gamma=\frac{r^2+H^2+2rHH_r}{(r^2-H^2)^2}.
 \label{model:definitions}
\end{equation}
Since \(P_r=R\) and \(Q_r=S\), we have
\begin{equation}
 G_r=PS_r-QR_r,\qquad
 \left(\frac RG\right)_r
 =\frac{R_rG-RG_r}{G^2}
 =-\frac{P\mathcal W}{G^2}.
 \label{model:quotient-derivative}
\end{equation}
It therefore suffices to prove \(\mathcal W\geq0\).

For a quadratic profile, raising the base height while keeping the peak
fixed changes the base slope at rate \(\kappa\).  The corresponding
transport operator on model data is
\begin{equation}
 \mathcal D=\kappa\partial_\lambda-\alpha\partial_r.
 \label{model:transport-operator}
\end{equation}
The arc transport identity from Section~\ref{sec:arcs} becomes
\(\mathcal DH=H/(r^2-H^2)-\lambda\).
Differentiating this identity in \(\lambda\) and in \(\kappa\), including
the derivatives of \(\mathcal D\), gives
\begin{equation}
 \mathcal DP=\gamma P-1,\qquad
 \mathcal DQ=\gamma Q-P.
 \label{model:PQ-transport}
\end{equation}
For example, differentiating in \(\lambda\) contributes
\(-\alpha_\lambda H_r\) on the left, where
\(\alpha_\lambda=2rHP/(r^2-H^2)^2\).  Differentiating in \(\kappa\)
also contributes \(P\), because the coefficient of
\(\partial_\lambda\) is \(\kappa\).  These are exactly the terms
accounted for in Equation~\eqref{model:PQ-transport}.

The commutator identity is
\(\mathcal D(f_r)=(\mathcal Df)_r+\alpha_r f_r\).
It yields
\begin{align}
 \mathcal DR&=(\gamma+\alpha_r)R+\gamma_rP,\nonumber\\
 \mathcal DS&=(\gamma+\alpha_r)S+\gamma_rQ-R,
 \label{model:RS-transport}\\
 \mathcal D(R_r)&=(\gamma+2\alpha_r)R_r
                  +(2\gamma_r+\alpha_{rr})R+\gamma_{rr}P,\nonumber\\
 \mathcal D(S_r)&=(\gamma+2\alpha_r)S_r
                  +(2\gamma_r+\alpha_{rr})S+\gamma_{rr}Q-R_r.
 \label{model:RSr-transport}
\end{align}
Applying these four formulas to \(RS_r-SR_r\), the terms containing
\(\alpha_{rr}\) cancel, and the remaining terms give
\begin{equation}
 \mathcal D\mathcal W
 =(2\gamma+3\alpha_r)\mathcal W
   +\gamma_rG_r-\gamma_{rr}G.
 \label{model:W-transport}
\end{equation}

The variation formula in Lemma~\ref{lem:variation}, together with the
smooth arc scaling in Section~\ref{sec:arcs}, gives the following
differentiated small-width expansions:
\begin{align}
 P&=\frac{r^2}{3}+O(r^3),&
 R&=\frac{2r}{3}+O(r^2),&
 R_r&=\frac23+O(r),\nonumber\\
 Q&=\frac{r^4}{15}+O(r^5),&
 S&=\frac{4r^3}{15}+O(r^4),&
 S_r&=\frac{4r^2}{5}+O(r^3).
 \label{model:small-width}
\end{align}
The remainders are uniform on compact sets of finite
\((\lambda,\kappa)\).  Indeed, after width-one scaling, the model
parameters are \(\lambda r,\kappa r^3\), the limiting arc is a regular
parabola, and the first variations in the profiles \(u\) and \(u^2/2\)
are \(1/3\) and \(1/15\).  Smooth dependence at that scaled arc permits
the two displayed \(r\)-differentiations.  In particular,
\begin{equation}
 \mathcal W=\frac{16}{45}r^3+O(r^4).
 \label{model:W-small-width}
\end{equation}

At a zero of \(\mathcal W\), the already established inequality \(R>0\)
allows us to set \(\sigma_0=R_r/R\).  Then
\begin{equation}
 R_r=\sigma_0R,\qquad S_r=\sigma_0S,\qquad G_r=\sigma_0G.
 \label{model:zero-relations}
\end{equation}
At such a zero, Equation~\eqref{model:W-transport} becomes
\[
 \mathcal D\mathcal W
 =-G(\gamma_{rr}-\sigma_0\gamma_r).
\]
The outward derivative along a characteristic is
$-\mathcal D/\alpha$. Thus the sign needed to prevent a first zero
from positive values is
\begin{equation}
 \gamma_{rr}-\sigma_0\gamma_r>0
 \qquad(\kappa>0,\ \mathcal W=0).
 \label{model:required-zero-sign}
\end{equation}
The next three subsections establish this strict inequality from the
quadratic arc equation. We then apply it on the entire model domain.

\subsection{Endpoint identities and two exact algebraic reductions}

For the remainder of the sign calculation, write
\(\kappa=2b>0\) and \(d=\lambda\).  The profile is \(du+bu^2\).
Its endpoints at base zero are \(-m,n\), where
\begin{equation}
 m=r-H,\qquad n=r+H,\qquad
 \ell=\frac{dn}{dm}=\frac{1+H_r}{1-H_r}>0,\qquad
 w=\frac{\ell_m}{\ell}.
 \label{model:endpoint-variables}
\end{equation}
Here \(m_r=1-H_r>0\), so \(m\) is a valid local coordinate along the
fixed-\((d,b)\) family.  Subscripts \(m\) below refer to this coordinate.
We also write
\begin{equation}
 \mathcal S=\frac{\ell_{mm}}{\ell}-\frac32w^2,
 \qquad
 w_m=\mathcal S+\frac12w^2.
 \label{model:Schwarzian-definition}
\end{equation}
Thus \(\mathcal S\) is the Schwarzian derivative of \(n=n(m)\).

The following weighted endpoint quantities will be useful:
\begin{align}
 X&=\frac{\ell/m-1/n}{1+\ell}=\kappa P+d,\nonumber\\
 Z&=\frac{\ell m-n}{1+\ell}=dP+3\kappa Q,\nonumber\\
 \gamma&=\frac{\ell/m^2+1/n^2}{1+\ell}.
 \label{model:endpoint-identities}
\end{align}
For the first identity, use Equation~\eqref{eq:fit-width-identity},
\[
 \kappa P+d=\frac{rH_r+H}{r^2-H^2},
\]
and substitute \(r=(m+n)/2\), \(H=(n-m)/2\).
For the second, width scaling gives
\begin{equation}
 H(\lambda,\kappa,r)
 =rH(\lambda r,\kappa r^3,1),\qquad
 rH_r-H=\lambda P+3\kappa Q.
 \label{model:scaling}
\end{equation}
The endpoint form of \(rH_r-H\) is the displayed expression for \(Z\).
The formula for \(\gamma\) follows directly from
Equation~\eqref{model:definitions}.

Define
\begin{align}
 D_*={}&(m-n)\mathcal S
 +\left[1-\frac{2n}{m}
       +\ell\left(\frac{2m}{n}-1\right)\right]w
 +\frac{m+n}{2}w^2,
 \label{model:Dstar}\\
 w_*={}&
 \frac{2(\ell m^4-2\ell m^3n+2mn^3-n^4)}
      {m^2n^2(m+n)}.
 \label{model:wstar}
\end{align}

\begin{lemma}
\label{model:endpoint-algebra}
The exact determinant identity is
\begin{equation}
 X_mZ_{mm}-Z_mX_{mm}
 =\frac{\ell^2(m+n)^2}{m^2n^2(1+\ell)^3}D_*
 =\frac{3\kappa^2}{m_r^3}\mathcal W.
 \label{model:determinant-factor}
\end{equation}
At a zero of \(\mathcal W\) with \(m\ne n\), set
\(\sigma=X_{mm}/X_m\).  This is well defined, and
\begin{align}
 \gamma_{rr}-\sigma_0\gamma_r
 &=m_r^2(\gamma_{mm}-\sigma\gamma_m),
 \label{model:coordinate-zero-sign}\\
 \gamma_{mm}-\sigma\gamma_m
 &=\frac{\ell(m+n)^2}{(1+\ell)m^2n^2}
       \frac{w_*-w}{m-n}.
 \label{model:gamma-factor}
\end{align}
All the prefactors displayed in
Equations~\eqref{model:determinant-factor} and
\eqref{model:gamma-factor} are strictly positive.
\end{lemma}

\begin{proof}
Put \(B=1/n^2-1/m^2\).  Differentiation of
Equation~\eqref{model:endpoint-identities}, using \(n_m=\ell\), gives
\begin{equation}
 Z_m=\frac{w\ell(m+n)}{(1+\ell)^2},\qquad
 X_m=\frac{\ell B}{1+\ell}+\frac{Z_m}{mn},\qquad
 B_m=\frac2{m^3}-\frac{2\ell}{n^3}.
 \label{model:first-endpoint-derivatives}
\end{equation}
Differentiate the middle identity and insert the first one into the
determinant.  After division by \(\ell^2(m+n)/(1+\ell)^3\), the result is
\begin{equation}
 B\mathcal S+
 \left(\frac{B(1+\ell)}{m+n}-B_m\right)w
 +\frac{(m+n)^2}{2m^2n^2}w^2.
 \label{model:determinant-intermediate}
\end{equation}
For example, in this calculation
\[
 \frac{d}{dm}\log\frac{\ell(m+n)}{(1+\ell)^2}
 =\frac{1-\ell}{1+\ell}w+\frac{1+\ell}{m+n},
 \qquad
 \left(\frac1{mn}\right)_m=-\frac{n+m\ell}{m^2n^2}.
\]
These two derivatives, together with
\(w_m=\mathcal S+w^2/2\), give the three coefficients in
Equation~\eqref{model:determinant-intermediate}.
Now
\[
 B=\frac{(m-n)(m+n)}{m^2n^2},
\]
and substitution of \(B_m\) shows that the coefficient of \(w\) in
Equation~\eqref{model:determinant-intermediate} is
\[
 \frac{m+n}{m^2n^2}
 \left[1-\frac{2n}{m}
       +\ell\left(\frac{2m}{n}-1\right)\right].
\]
This proves its first factorization in
Equation~\eqref{model:determinant-factor}.

On the other hand, differentiating \(X=\kappa P+d\) and
\(Z=dP+3\kappa Q\) with respect to \(r\) gives
\[
 X_rZ_{rr}-Z_rX_{rr}=3\kappa^2(RS_r-SR_r).
\]
The chain rule gives
\[
 X_rZ_{rr}-Z_rX_{rr}
 =m_r^3\bigl(X_mZ_{mm}-Z_mX_{mm}\bigr).
\]
This proves the other equality in
Equation~\eqref{model:determinant-factor}.  It also gives
\begin{equation}
 X_m=\frac{\kappa R}{m_r}>0.
 \label{model:Xm-positive}
\end{equation}
Thus, at a zero, \(\sigma=X_{mm}/X_m\) is legitimate, even if
\(Z_m=0\) or \(w=0\).  The determinant identity then implies
\(Z_{mm}=\sigma Z_m\).
Using \(X_{rr}=\sigma_0X_r\) at that point gives
\[
 \sigma m_r^2=\sigma_0m_r-m_{rr}.
\]
The chain rule for \(\gamma\) proves
Equation~\eqref{model:coordinate-zero-sign}.

It remains to compute \(\gamma_{mm}-\sigma\gamma_m\).
For a function \(f\) smooth near \(m\) and \(-n\), write
\[
 T_f=\frac{\ell f(m)+f(-n)}{1+\ell},\qquad
 a_0=\frac{\ell}{1+\ell},\qquad
 \beta=\frac{w}{1+\ell},\qquad
 \rho=(\log a_0)_m-\sigma.
\]
A first differentiation gives
\[
 (T_f)_m=a_0\{f'(m)-f'(-n)+\beta[f(m)-f(-n)]\}.
\]
A second differentiation therefore gives
\begin{align}
 \frac{(T_f)_{mm}-\sigma(T_f)_m}{a_0}
 ={}&f''(m)+\ell f''(-n)
       +\beta\{f'(m)+\ell f'(-n)\}\nonumber\\
 &+\rho\{f'(m)-f'(-n)\}
       +(\beta_m+\rho\beta)\{f(m)-f(-n)\}.
 \label{model:weighted-operator}
\end{align}
At the zero under consideration, this expression vanishes for
\(f(x)=x\) and \(f(x)=1/x\), since \(T_x=Z\) and \(T_{1/x}=X\).
Consequently
\begin{align}
 \beta_m+\rho\beta
   &=-\frac{\beta(1+\ell)}{m+n},\nonumber\\
 \rho B
   &=\frac{2\ell}{n^3}-\frac2{m^3}
     +\beta\left(\frac1{m^2}+\frac{\ell}{n^2}
                         +\frac{1+\ell}{mn}\right).
 \label{model:rho-relations}
\end{align}
Here \(B\ne0\) because \(m\ne n\); no division by \(\beta\) is made.

For \(f(x)=1/x^2\), Equation~\eqref{model:weighted-operator} becomes
\begin{align}
 \frac{\gamma_{mm}-\sigma\gamma_m}{a_0}
 ={}&\frac6{m^4}+\frac{6\ell}{n^4}
       +\beta\left(-\frac2{m^3}+\frac{2\ell}{n^3}\right)
       -2\rho\left(\frac1{m^3}+\frac1{n^3}\right)
       +\frac{\beta(1+\ell)B}{m+n}.
 \label{model:gamma-intermediate}
\end{align}
Eliminate \(\rho\) using Equation~\eqref{model:rho-relations}.
The resulting constant coefficient and coefficient of \(\beta\) are,
respectively,
\begin{align}
 A_0={}&\frac6{m^4}+\frac{6\ell}{n^4}
   -\frac{2(1/m^3+1/n^3)(2\ell/n^3-2/m^3)}{B}
 =\frac{2(m+n)N}{m^4n^4(m-n)},\nonumber\\
 A_\beta={}&-\frac2{m^3}+\frac{2\ell}{n^3}
 -\frac{2(1/m^3+1/n^3)
       (1/m^2+\ell/n^2+(1+\ell)/(mn))}{B}
 +\frac{(1+\ell)B}{m+n}\nonumber\\
 ={}&-\frac{(1+\ell)(m+n)^2}{m^2n^2(m-n)},
 \label{model:gamma-coefficients}\\
 N={}&\ell m^4-2\ell m^3n+2mn^3-n^4.\nonumber
\end{align}
For completeness, the two numerator reductions in
Equation~\eqref{model:gamma-coefficients} are
\begin{align*}
 &6(n^4+\ell m^4)(m^2-n^2)
       -4(m^3+n^3)(\ell m^3-n^3)=2(m+n)^2N,\\
 &(m+n)\bigl[2(\ell m^3-n^3)(m-n)
       -2(m^3+n^3)(\ell m+n)
       +(1+\ell)mn(m-n)^2\bigr]\\
 &\hspace{35mm}=-(1+\ell)mn(m+n)^3.
\end{align*}
They correspond to the common denominators
\(m^4n^4(m^2-n^2)\) and \(m^3n^3(m^2-n^2)\).
Substituting \(\beta=w/(1+\ell)\) into \(A_0+\beta A_\beta\)
now gives
\[
 \frac{(m+n)^2}{m^2n^2}\frac{w_*-w}{m-n}.
\]
Multiplication by \(a_0\) proves
Equation~\eqref{model:gamma-factor}.
\end{proof}

\subsection{Riccati equations, a common endpoint label, and an estimate}

The algebra has reduced the desired sign at a possible zero to the
ordering of $w$ and $w_*$. To establish that ordering we now use
the differential equation of the actual quadratic arc. Its Riccati
linearization gives a common projective label for the two endpoints;
the classical Schwarzian ratio identity and chain rule
\cite{OT} then determine the endpoint Schwarzian.
The next lemma derives these relations and the additional integral
estimate used in the sign argument.

\begin{lemma}
\label{model:Riccati}
For the quadratic profile \(du+bu^2\), with \(b>0\), one has
\begin{equation}
 \mathcal S
 =2b(m+n\ell^2)+\frac{d^2}{2}(\ell^2-1)
   -d\left(\frac1m+\frac{\ell^2}{n}\right)
   +\frac32\left(\frac{\ell^2}{n^2}-\frac1{m^2}\right).
 \label{model:Schwarzian-formula}
\end{equation}
The arc admits a positive function \(V\) such that
\begin{align}
 \ell&=\frac{mV(-m)^2}{nV(n)^2},\nonumber\\
 w&=\frac1m-\frac{\ell}{n}+d(1+\ell)
       +2bmV(-m)^2\int_{-m}^{n}V(y)^{-2}\,dy.
 \label{model:variation-formulas}
\end{align}
If \(d<0\) and \(d^2\ge4bm\), put
\(\nu=(d^2/4-bm)^{1/2}\).  Then
\begin{equation}
 w\le\frac1m-\frac{\ell}{n}+d\ell-2\nu.
 \label{model:Riccati-bound}
\end{equation}
This includes \(\nu=0\).
\end{lemma}

\begin{proof}
The arc equation is \(u_y=du+bu^2-y\).  On its finite interval
\([-m,n]\), define
\begin{equation}
 V(y)=\exp\left[-\int_{-m}^{y}(bu(s)+d/2)\,ds\right]>0,
 \qquad U(y)=u(y)V(y).
 \label{model:positive-V}
\end{equation}
The integral is finite because the actual arc is bounded on a compact
interval. In the following system, primes denote derivatives with
respect to \(y\). Direct differentiation gives the regular linear system
\begin{equation}
 U'=\frac d2U-yV,\qquad V'=-bU-\frac d2V.
 \label{model:linear-system}
\end{equation}
In particular \(U(-m)=U(n)=0\).
Eliminating \(V\) where \(y\ne0\) gives
\begin{equation}
 U''=\frac{U'}{y}
       +\left(by+\frac{d^2}{4}-\frac{d}{2y}\right)U.
 \label{model:scalar-U}
\end{equation}

We now justify the use of one projective solution label at both
endpoints despite the singularity of Equation~\eqref{model:scalar-U}
at zero.  Take \((U_1,V_1)=(U,V)\) from
Equation~\eqref{model:positive-V}, and complete it to a fundamental
pair \((U_1,V_1),(U_2,V_2)\) of the regular system in
Equation~\eqref{model:linear-system}.  Its coefficient matrix has trace zero, so
\[
 \Omega=U_1V_2-U_2V_1
\]
is a nonzero constant.  At either endpoint, \(U_1=0\) and \(V_1>0\);
hence \(U_2=-\Omega/V_1\ne0\).  Therefore the same ratio
\(f=U_1/U_2\) is finite near both endpoints, and
\begin{equation}
 f(-m)=f(n)=0,\qquad f'(y)=\frac{y\Omega}{U_2(y)^2}\ne0
 \quad\hbox{at }y=-m,n.
 \label{model:common-label}
\end{equation}
A nearby arc has coefficients that can be normalized as \((1,-q)\)
in this fixed fundamental pair: continuous dependence keeps its first
coefficient nonzero near the original arc.  Its endpoint condition is
\(U_1-qU_2=0\).  Thus the nearby endpoint functions satisfy the same
local identity \(f(n(m))=f(-m)\).  No scalar solution is continued
through the singular equation at zero.  The scalar determinant
\(U_1'U_2-U_1U_2'=y\Omega\) is nonzero in each endpoint neighborhood, so \(U_1,U_2\) form
a scalar basis in each such neighborhood.

Writing \(U=\sqrt{|y|}\,\widetilde U\) in
Equation~\eqref{model:scalar-U} gives
\[
 \widetilde U''
 =\left(by+\frac{d^2}{4}-\frac{d}{2y}
                    +\frac3{4y^2}\right)\widetilde U.
\]
For two solutions of \(a''=Q_0a\), \(c''=Q_0c\), their ratio has
Schwarzian \(-2Q_0\): locally,
\((a/c)''/(a/c)'=-2c'/c\), and differentiating this expression and
subtracting half its square gives \(-2c''/c\).
The ratio is unchanged by division of both solutions by
\(\sqrt{|y|}\).  Consequently
\[
 \{f,y\}=-2by-\frac{d^2}{2}+\frac d y-\frac3{2y^2},
 \qquad
 \{f,y\}=\frac{f'''}{f'}-\frac32\left(\frac{f''}{f'}\right)^2.
\]
Apply the Schwarzian chain rule to the common-label identity
\(f(n(m))=f(-m)\).  The affine map \(m\mapsto-m\) has zero
Schwarzian and squared derivative one, so
\[
 \mathcal S=\{f,y\}\big|_{y=-m}
              -\ell^2\{f,y\}\big|_{y=n}.
\]
Expanding these two evaluations proves
Equation~\eqref{model:Schwarzian-formula}.

For Equation~\eqref{model:variation-formulas}, vary the initial endpoint
in \(u(-m;m)=0\).  Since \(u_y(-m)=m\), the initial variation is
\(u_m(-m)=m\).  Its scalar variational equation yields
\[
 u_m(y)=m\exp\left[\int_{-m}^{y}(d+2bu(s))\,ds\right]
       =m\,\frac{V(-m)^2}{V(y)^2}.
\]
Differentiating \(u(n(m);m)=0\), and using \(u_y(n)=-n\), gives the
formula for \(\ell\).  Its logarithm can be written
\[
 \log\ell=\log m-\log n+\int_{-m}^{n}(d+2bu(y))\,dy.
\]
Differentiate in \(m\).  The two boundary contributions from \(2bu\)
vanish, while the constant term contributes \(d(1+\ell)\).
Substituting the expression for \(u_m\) proves the formula for \(w\).

Finally, differentiating the second equation in
Equation~\eqref{model:linear-system} gives
\[
 V''=(by+d^2/4)V,\qquad V'(-m)=-dV(-m)/2.
\]
Suppose \(d<0\), \(d^2\ge4bm\), and set \(s=y+m\).
For \(v_0(s)=V(-m+s)/V(-m)\) one has
\[
 v_0''=(\nu^2+bs)v_0,\qquad v_0(0)=1,\qquad v_0'(0)=-d/2.
\]
Let
\[
 f_\nu(s)=
 \begin{cases}
 \displaystyle\cosh(\nu s)-\frac d{2\nu}\sinh(\nu s),&\nu>0,\\[1mm]
 1-ds/2,&\nu=0.
 \end{cases}
\]
Variation of constants gives, for \(0\le s\le m+n\),
\[
 v_0(s)-f_\nu(s)
 =b\int_0^s K_\nu(s-t)\,t\,v_0(t)\,dt\ge0,
 \qquad
 K_\nu(t)=
 \begin{cases}\sinh(\nu t)/\nu,&\nu>0,\\t,&\nu=0.\end{cases}
\]
The inequality follows from the already constructed positivity of
\(v_0\), and from \(K_\nu\ge0\).  Also \(f_\nu>0\) on the positive
half-line because \(d<0\).  Therefore
\begin{equation}
 V(-m)^2\int_{-m}^{n}V(y)^{-2}\,dy
 \le\int_0^\infty f_\nu(s)^{-2}\,ds
 =\frac1{\nu-d/2}.
 \label{model:inverse-square-bound}
\end{equation}
For \(\nu>0\), the substitution \(q=\tanh(\nu s)\), with
\(A=-d/(2\nu)>0\), evaluates the last integral as
\[
 \frac1\nu\int_0^1(1+Aq)^{-2}\,dq
 =\frac1{\nu(1+A)}.
\]
For \(\nu=0\), direct integration of \((1-ds/2)^{-2}\) gives
\(-2/d\), the same formula.
Since \(\nu^2=d^2/4-bm\),
\[
 \frac{2bm}{\nu-d/2}=-d-2\nu.
\]
Substitution into Equation~\eqref{model:variation-formulas} proves
Equation~\eqref{model:Riccati-bound}.
\end{proof}

\subsection{The signs at a possible zero}

\begin{lemma}
\label{model:zero-sign}
If \(\kappa>0\) and \(\mathcal W=0\), then
\(\gamma_{rr}-\sigma_0\gamma_r>0\), with \(\sigma_0\) defined in
Equation~\eqref{model:zero-relations}.
\end{lemma}

\begin{proof}
We partition the possibilities for \(m,n,\ell\).
The quantities \(\theta=\operatorname{arctanh}(H/r)\) and
\(L=\operatorname{arctanh}(H_r)\), and the finite-peak reference
characteristics used below, are those of Lemma~\ref{fit:references}.

\paragraph{Case 1: \(m\le n\).}
The final angle is nonnegative.  Along the characteristic leading to
this data point, the driving slope satisfies
\[
 j'=-s\kappa\operatorname{sech}^2\theta<0.
\]
Its initial value \(z\) must be positive.  Indeed, if \(z\le0\), then
\(j(s)<0\) for every \(s>0\), and
\(\theta'+a(s)\theta=j(s)\), where
\(a(s)=\sinh(2\theta)/(s\theta)>0\), with value \(2/s\) at
\(\theta=0\), forces \(\theta(s)<0\).  This last assertion follows by
integrating from a positive \(\varepsilon\) and letting
\(\varepsilon\downarrow0\): \(\theta(\varepsilon)\to0\), and the
forward homogeneous factor is at most one.
It would contradict the nonnegative final angle.

Since \(z>0\), \(\theta(s)=zs/3+O(s^2)>0\) initially.
If it first reached zero at some earlier width \(s_1<r\), then
\(j(s_1)=\theta'(s_1)\le0\).  The strict decrease of \(j\) would give
\(j(s)<0\) thereafter, forcing \(\theta(s)<0\) for \(s>s_1\).
Again this contradicts the final angle.  Thus
\(\theta(s)>0\) for all \(0<s<r\), and
\[
 L(r)=\int_0^r\frac{\sinh(2\theta(s))}{s}\,ds>0,
 \qquad \ell=e^{2L}>1.
\]

We also need the sign of \(w=(\log\ell)_m\), where the derivative
is taken along the fixed-profile family at base zero. Let \(t=t(r)\)
be the peak height of that profile.
Lemma~\ref{lem:arc-data} gives \(t_r>0\), and its peak slope is
\(z=d+\kappa t\).  The final model value \(L=\operatorname{arctanh}H_r\)
equals \(\mathcal L_\kappa(r,z(r))\), and
Lemma~\ref{fit:references} gives \(\partial_z\mathcal L_\kappa>0\).
Consequently
\begin{equation}
 L_r\big|_{d,\kappa}
 =\frac{\sinh(2\theta)}r
   +\kappa t_r\partial_z\mathcal L_\kappa(r,z)>0,\qquad
 w=\frac{2L_r|_{d,\kappa}}{m_r}>0.
 \label{model:w-positive-nonnegative-angle}
\end{equation}
If \(m=n\), Equation~\eqref{model:Dstar} now reduces to
\[
 D_*=(\ell-1)w+mw^2>0.
\]
Hence a zero of \(\mathcal W\) cannot have \(m=n\).

If \(m<n\), put \(z_1=m/n\in(0,1)\).  The sign of \(w_*\) is the sign
of
\[
 \ell z_1^3(z_1-2)+2z_1-1
 \le z_1^3(z_1-2)+2z_1-1
 =(z_1-1)^3(z_1+1)<0.
\]
Since \(w>0\), the quotient \((w_*-w)/(m-n)\) is positive.
Equations~\eqref{model:gamma-factor} and
\eqref{model:coordinate-zero-sign} prove the required sign.

\paragraph{Case 2: \(m>n\) and \(\ell\ge1\).}
Here the final angle is negative, which forces \(d<0\).
In fact, if its final driving slope \(d\) were nonnegative, then the
strictly decreasing driving slope would be positive at all earlier
widths, and the same scalar angle equation would give a positive final
angle.

Each group in Equation~\eqref{model:Schwarzian-formula} is now
nonnegative, and several are strictly positive.  In particular
\begin{equation}
 \mathcal S>0
 \quad\hbox{whenever }m>n,\ \ell\ge1.
 \label{model:S-positive-ell-large}
\end{equation}
We justify \(w\ge0\) on this region without assuming that \(m\)
ranges to infinity.

For the fixed quadratic profile, Lemma~\ref{lem:arc-data} gives arcs at
all widths \(r\in(0,\infty)\).  Since \(m_r>0\) and \(0<m(r)<2r\),
the \(m\)-domain is an interval \((0,M_*)\), with \(M_*\) finite or
infinite.  At its left end,
\[
 \ell=1+\frac43dr+O(r^2)<1
\]
because \(d<0\).
Fix a point in the open set \(\{m-n>0,\ell>1\}\), and let \((a,c)\)
be its connected component.  The expansion just given implies \(a>0\);
also \(a<M_*\), since the chosen point is to its right.
Thus \(a\) is an interior point of the smooth parameter domain.
On the component,
\((m-n)_m=1-\ell<0\).  Moving backwards to \(a\) therefore increases
\(m-n\), so its value at \(a\) is strictly positive.
The component cannot enter through \(m=n\).  It must have
\(\ell(a)=1\), and the right derivative there satisfies
\(\ell_m(a)\ge0\), or \(w(a)\ge0\).
Equations~\eqref{model:Schwarzian-definition} and
\eqref{model:S-positive-ell-large} give \(w_m>0\) throughout the
component.  Hence \(w>0\) inside it.

At a point with \(m>n,\ell=1,w<0\), points immediately to the left
would have \(\ell>1,w<0\), contrary to the result just proved.
A point with \(m>n,\ell=1,w=0\) is also impossible: there
\(w_m=\mathcal S>0\) and
\(\ell_{mm}=\ell(w_m+w^2)=\mathcal S>0\).
Immediately to the left one would again have \(\ell>1,w<0\).
Thus \(w\ge0\) at every point required in this case.
The component argument applies to each chosen component separately;
it makes no assumption about their number.

Writing \(z_1=m/n>1\), the coefficient of \(w\) in
Equation~\eqref{model:Dstar} satisfies
\[
 1-\frac2{z_1}+\ell(2z_1-1)
 \ge 2z_1-\frac2{z_1}>0.
\]
Together with \((m-n)\mathcal S>0\) and \(w\ge0\), this gives
\(D_*>0\).  Hence no zero of \(\mathcal W\) occurs in this case.

\paragraph{Case 3: \(m>n\) and \(0<\ell<1\).}
Again \(d<0\).  It suffices to prove that
\begin{equation}
 w\ge w_*\quad\Longrightarrow\quad D_*>0.
 \label{model:last-case-target}
\end{equation}
Indeed, at a zero this will give \(w<w_*\), and the denominator
\(m-n\) in Equation~\eqref{model:gamma-factor} is positive.

We may set \(n=1,m=z_1>1\) at the point.  To specify this scaling, fix
the original endpoint value \(N=n>0\) as a constant and put
\(y=N\widetilde y\), \(u=N^2\widetilde u\).
The new model parameters are
\(\widetilde d=Nd\), \(\widetilde b=N^3b\), while
\[
 \widetilde m=m/N,\quad \widetilde n=n/N,\quad
 \widetilde\ell=\ell,\quad
 \widetilde w=Nw,\quad
 \widetilde{\mathcal S}=N^2\mathcal S,\quad
 \widetilde D_*=ND_*,\quad \widetilde w_*=Nw_*.
\]
Thus the equations and the signs needed for
Equation~\eqref{model:last-case-target} are preserved.  We suppress
the tildes and use \(n=1,m=z_1\) from now on.

In these units,
\begin{equation}
 w_*=\frac{2(\ell z_1^3(z_1-2)+2z_1-1)}{z_1^2(z_1+1)}>0.
 \label{model:scaled-wstar}
\end{equation}
For \(z_1\ge2\) positivity is immediate.  For \(1<z_1<2\), the coefficient
of \(\ell\) in its numerator is negative; replacing \(\ell<1\)
by one decreases that numerator to
\((z_1-1)^3(z_1+1)>0\).

For fixed \(\mathcal S\), write
\[
 D_*=(z_1-1)\mathcal S+
       [1-2/z_1+\ell(2z_1-1)]w+\frac{z_1+1}{2}w^2.
\]
The derivative of this expression in \(w\), evaluated at \(w_*\) and
multiplied by \(z_1^2\), is
\begin{equation}
 J(\ell,z_1)=2\ell z_1^4-2\ell z_1^3-\ell z_1^2+z_1^2+2z_1-2.
 \label{model:J-polynomial}
\end{equation}
This polynomial is affine in \(\ell\).  With \(\xi=z_1-1>0\), its
endpoint values are
\[
 J(0,1+\xi)=\xi^2+4\xi+1,\qquad
 J(1,1+\xi)=2\xi(\xi+2)(\xi^2+\xi+1).
\]
Both are strictly positive.  Therefore \(J>0\) for \(0<\ell<1\).
Since the derivative in \(w\) increases with \(w\), \(D_*\) is
strictly increasing for \(w\ge w_*\).

Set
\begin{equation}
 T=2\ell z_1^3-\ell z_1^2-z_1+2,\qquad
 A=2\ell z_1^3-3\ell z_1^2+3z_1-2,\qquad
 \mathcal S_0=\frac32(\ell^2-z_1^{-2}).
 \label{model:T-and-baseline}
\end{equation}
Substitution of Equation~\eqref{model:scaled-wstar} gives the two
polynomial identities
\begin{align}
 (z_1-1)\mathcal S_0+
 [1-2/z_1+\ell(2z_1-1)]w_*+\frac{z_1+1}{2}w_*^2
 &=\frac{z_1-1}{2z_1^4}AT,
 \label{model:baseline-factorization}\\
 \Delta_0:=\frac1{z_1}-\ell-w_*
 &=-\frac{z_1-1}{(z_1+1)z_1^2}T.
 \label{model:Delta-factorization}
\end{align}
These factorizations can be checked by collecting powers of \(\ell\).
For Equation~\eqref{model:baseline-factorization}, multiplying its left
side by \(2z_1^4(z_1+1)\) gives
\[
 (z_1^2-1)\bigl[
   \ell^2z_1^4(2z_1-3)(2z_1-1)
   +4\ell z_1^2(z_1^2-1)+(3z_1-2)(2-z_1)\bigr].
\]
The bracket is \(AT\), because the cross coefficient in the product
is
\[
 z_1^2\bigl[(2z_1-3)(2-z_1)+(3z_1-2)(2z_1-1)\bigr]
 =4z_1^2(z_1^2-1).
\]
For Equation~\eqref{model:Delta-factorization}, the numerator over
\(z_1^2(z_1+1)\) is
\((z_1-1)[z_1-2-\ell z_1^2(2z_1-1)]=-(z_1-1)T\).
The sign of \(A\) requires no further restriction on the parameters.
It too is affine in \(\ell\), and
\[
 A(0,1+\xi)=3\xi+1,\qquad
 A(1,1+\xi)=\xi(2\xi^2+3\xi+3),
\]
so \(A>0\).

Assume \(w\ge w_*\). First suppose \(T\ge0\).
If \(d^2\ge4bz_1\), Equation~\eqref{model:Riccati-bound}
would imply
\[
 \Delta_0
 =\frac1{z_1}-\ell-w_*
 \ge\frac1{z_1}-\ell-w
 \ge-d\ell+2\nu>2\nu\ge0.
\]
This contradicts Equation~\eqref{model:Delta-factorization}.
Consequently \(d^2<4bz_1\).  In Equation~\eqref{model:Schwarzian-formula},
the difference from the baseline is
\begin{equation}
 \mathcal S-\mathcal S_0
 =\left(2bz_1-\frac{d^2}{2}\right)
   +\ell^2\left(2b+\frac{d^2}{2}-d\right)-\frac d{z_1}>0.
 \label{model:baseline-improvement}
\end{equation}
Each term on the right is positive in this subcase.
Equation~\eqref{model:baseline-factorization} is nonnegative because
\(A>0,T\ge0\).  Since \(z_1-1>0\), the strict improvement in
Equation~\eqref{model:baseline-improvement}, followed by monotonicity
in \(w\), gives \(D_*>0\).

It remains to consider \(T<0\).  Then
\[
 z_1-2>\ell z_1^2(2z_1-1)>0.
\]
In particular \(z_1>2\), and the same inequality implies \(z_1-2>3\ell\).
Equation~\eqref{model:endpoint-identities}, with \(\kappa P>0\), gives
\[
 d<\frac{\ell/z_1-1}{1+\ell},\qquad
 -dz_1>\frac{z_1-\ell}{1+\ell}>2.
\]
Rewrite Equation~\eqref{model:Schwarzian-formula} as
\begin{equation}
 \mathcal S=
 \left(2bz_1-\frac{d^2}{2}\right)
 +\ell^2\left(2b+\frac{d^2}{2}-d+\frac32\right)
 -\frac d{z_1}-\frac3{2z_1^2}.
 \label{model:S-last-subcase}
\end{equation}
If the first parentheses are nonnegative, then
\(-d/z_1>2/z_1^2\) immediately gives \(\mathcal S>0\).
If they are negative, put \(\nu=(d^2/4-bz_1)^{1/2}>0\).
The preceding use of Equation~\eqref{model:Riccati-bound} still gives
\(2\nu<\Delta_0\), while \(w_*>0,\ell>0\) give
\(\Delta_0<1/z_1\).  Hence
\[
 2bz_1-\frac{d^2}{2}=-2\nu^2>-\frac1{2z_1^2}.
\]
Together with \(-d/z_1>2/z_1^2\), this strictly dominates the negative
term \(-3/(2z_1^2)\) in Equation~\eqref{model:S-last-subcase}.
Its remaining \(\ell^2\)-term is positive, so again
\(\mathcal S>0\).
Since \(z_1>2\), the coefficient \(1-2/z_1+\ell(2z_1-1)\) is positive,
and \(w\ge w_*>0\).  Therefore every term in \(D_*\) is positive.
This proves Equation~\eqref{model:last-case-target} also when \(T<0\).

The three cases cover all endpoint configurations.
At a zero of \(\mathcal W\), Cases 1 and 3 give
\((w_*-w)/(m-n)>0\), and the excluded diagonal and Case 2 contain no
zeros.  Equations~\eqref{model:gamma-factor} and
\eqref{model:coordinate-zero-sign} finish the proof.
\end{proof}

\subsection{Completion of the characteristic argument}

\begin{proof}[Proof of Theorem~\ref{thm:model-inequality}]
First suppose \(\kappa>0\).  Fix any model data point
\((\lambda,\kappa,r)\).  By the arc construction and
Lemma~\ref{fit:references}, it belongs to a characteristic coming
from a finite peak.  Along that characteristic, use \(s\) for
outward-increasing width.  The base slope tends to a finite peak slope
as \(s\downarrow0\), so the uniform expansion in
Equation~\eqref{model:W-small-width} makes \(\mathcal W\) strictly
positive for all sufficiently small positive \(s\).

On this curve \(\kappa\) is constant and
\(d\lambda/ds=-\kappa/\alpha\).  Thus the total outward derivative is
\begin{equation}
 \frac{d}{ds}\mathcal W(\lambda(s),\kappa,s)
 =\mathcal W_r-\frac{\kappa}{\alpha}\mathcal W_\lambda
 =-\frac1\alpha\,\mathcal D\mathcal W.
 \label{model:outward-orientation}
\end{equation}
At any zero, Equations~\eqref{model:W-transport} and
\eqref{model:zero-relations}, and Lemma~\ref{model:zero-sign}, give
\[
 \mathcal D\mathcal W
 =-G(\gamma_{rr}-\sigma_0\gamma_r)<0.
\]
Since \(\alpha>0\), Equation~\eqref{model:outward-orientation}
makes the outward derivative strictly positive.  A first zero reached
from strictly positive values would instead have a nonpositive
outward derivative.  This contradiction proves
\(\mathcal W>0\) at every point with \(\kappa>0\).

Reflection of the endpoint coordinate gives
\[
 H(-\lambda,-\kappa,r)=-H(\lambda,\kappa,r).
\]
Differentiating this symmetry shows that \(P,Q\), their
\(r\)-derivatives, and \(\mathcal W\) are unchanged by simultaneous
sign reversal of \(\lambda,\kappa\).
Thus \(\mathcal W>0\) also when \(\kappa<0\).
Smooth dependence of the model on finite parameters at every \(r>0\)
now gives \(\mathcal W\ge0\) at \(\kappa=0\).
Only this non-strict conclusion at zero curvature is required.
Locally the passage is uniform on compact parameter sets with \(r>0\);
\(G>0\) there is already known from
Theorem~\ref{thm:quadratic-fit}.

Finally, apply the pointwise fixed-parameter identity
Equation~\eqref{model:quotient-derivative}.
It gives \((R/G)_r\le0\) for every \(\lambda,\kappa,r\), as asserted.
The characteristic transport established the sign of a scalar
function on the whole model domain; it did not identify its partial
\(r\)-derivative with a total derivative along a characteristic.
\end{proof}
